\chapter{虚拟存储器}

\section{虚拟存储器概念}
\concept{为什么需要虚拟存储器}
\begin{points}
  \item 传统物理存储管理要求作业运行前全部装入内存，运行期间常驻内存，限制了大作业和多作业运行。
  \item 物理扩充内存会增加成本，因此需要从逻辑上扩充内存。
  \item 虚拟存储器打破“作业必须全部装入内存才能运行”的限制。
\end{points}
\plain{不是资源真的变多，而是 OS 用映射、换入换出、排队等办法，让用户感觉自己有一份独立资源。}
\exam{常考虚拟处理机、虚拟存储器、虚拟设备的例子，以及虚拟性和资源复用的关系。}


\concept{局部性原理}
\begin{points}
  \item 程序执行时，CPU 访问的指令和数据倾向于聚集在较小范围内。
  \item 时间局部性：刚访问过的指令或数据，短时间内可能再次访问。
  \item 空间局部性：访问某地址后，附近地址很可能被访问。
  \item 顺序局部性：程序大多数情况下顺序执行，下一条指令通常紧接当前指令。
  \item 局部性原理是虚拟存储器的理论基础。
\end{points}
\plain{程序一段时间只爱访问附近的一小块代码和数据，所以没必要一开始全装入内存。}
\exam{常考时间/空间局部性、虚拟存储器四特征、容量受地址结构和外存限制。}


\concept{虚拟存储器定义与特征}
\begin{points}
  \item 虚拟存储器是具有请求调入和置换功能、能从逻辑上扩充内存容量的存储器系统。
  \item 它是一种以时间换空间的技术：用缺页中断、磁盘换入换出换取更大的逻辑空间。
  \item 特征：离散性、多次性、对换性、虚拟性。
  \item 容量受地址结构和外存容量限制，不是无限大。
\end{points}
\plain{不是资源真的变多，而是 OS 用映射、换入换出、排队等办法，让用户感觉自己有一份独立资源。}
\exam{常考虚拟处理机、虚拟存储器、虚拟设备的例子，以及虚拟性和资源复用的关系。}


\section{请求分页存储管理}
\concept{请求分页的基本思想}
\begin{points}
  \item 程序运行前只装入部分页面即可启动。
  \item 运行中访问的页面若不在内存，产生缺页中断，由 OS 调入所需页面。
  \item 若内存无空闲页框，需要选择某个页面换出。
  \item 请求分页 = 基本分页 + 请求调页 + 页面置换。
\end{points}
\plain{分页就是逻辑上按页切、物理上按页框放，页表负责记录页到页框的映射。}
\exam{常考页号/页内偏移拆分、有效访问时间、多级页表节省页表空间。}


\concept{基本分页与请求分页区别}
\begin{points}
  \item 基本分页要求进程运行所需页面已在内存中，地址转换时页表项应有效。
  \item 请求分页允许部分页面不在内存，访问未驻留页面时产生缺页中断。
  \item 因此基本分页本身不以缺页中断为常规机制；请求分页才把缺页作为正常运行过程的一部分。
\end{points}
\plain{分页就是逻辑上按页切、物理上按页框放，页表负责记录页到页框的映射。}
\exam{常考页号/页内偏移拆分、有效访问时间、多级页表节省页表空间。}


\concept{页表项扩展}
\begin{points}
  \item 存在位/有效位：表示该页是否在内存。
  \item 修改位/脏位：表示页面是否被写过，换出时若为脏页需写回外存。
  \item 引用位/访问位：表示页面近期是否被访问，用于页面置换算法。
  \item 外存地址：指出页面在外存中的位置。
\end{points}
\plain{分页就是逻辑上按页切、物理上按页框放，页表负责记录页到页框的映射。}
\exam{常考页号/页内偏移拆分、有效访问时间、多级页表节省页表空间。}


\concept{缺页中断处理过程}
\begin{points}
  \item CPU 访问页面，发现页表项存在位为 0，产生缺页中断。
  \item 操作系统检查访问是否合法；若非法，终止进程。
  \item 若合法，查找空闲页框；若无空闲页框，执行页面置换。
  \item 若被换出页被修改过，需要写回外存。
  \item 从外存读入所需页面，更新页表和快表，重新执行引起缺页的指令。
\end{points}
\plain{中断就是硬件或软件打断 CPU 当前执行，让 OS 先处理更紧急或必须处理的事件。}
\exam{常考中断处理步骤、中断与异常/系统调用区别、为什么中断能提高 CPU 与 I/O 并行度。}


\section{页面置换算法}
\concept{最佳置换 OPT}
\begin{points}
  \item OPT 选择未来最长时间不会被访问的页面淘汰。
  \item 它缺页率最低，但需要预知未来，实际无法实现。
  \item 常作为评价其他算法的理论下界。
\end{points}
\plain{内存满了要踢页；OPT 看未来，FIFO 看入队时间，LRU 看最近使用，CLOCK 用引用位近似 LRU。}
\exam{常考给页面引用串和页框数计算缺页次数，FIFO 可能有 Belady 异常。}


\concept{先进先出 FIFO}
\begin{points}
  \item FIFO 淘汰最早进入内存的页面。
  \item 实现简单，只需维护队列。
  \item 缺点是不考虑页面使用频率，可能淘汰仍常用页面。
  \item FIFO 可能出现 Belady 异常：页框数增加，缺页次数反而增加。
\end{points}
\plain{内存满了要踢页；OPT 看未来，FIFO 看入队时间，LRU 看最近使用，CLOCK 用引用位近似 LRU。}
\exam{常考给页面引用串和页框数计算缺页次数，FIFO 可能有 Belady 异常。}


\concept{最近最久未使用 LRU}
\begin{points}
  \item LRU 淘汰最近最长时间没有被访问的页面。
  \item 基于局部性原理：最近很久未用的页面，未来短期内也可能不用。
  \item 实现可用计数器、栈或近似算法，但精确 LRU 开销较高。
  \item LRU 不会出现 Belady 异常。
\end{points}
\plain{内存满了要踢页；OPT 看未来，FIFO 看入队时间，LRU 看最近使用，CLOCK 用引用位近似 LRU。}
\exam{常考给页面引用串和页框数计算缺页次数，FIFO 可能有 Belady 异常。}


\concept{时钟 CLOCK/二次机会算法}
\begin{points}
  \item Clock 用访问位近似 LRU，把页面组成循环队列。
  \item 指针扫描页面：若访问位为 0，淘汰；若为 1，清 0 并给第二次机会。
  \item 改进型 Clock 还考虑修改位，优先淘汰未访问且未修改的页面。
\end{points}
\plain{内存满了要踢页；OPT 看未来，FIFO 看入队时间，LRU 看最近使用，CLOCK 用引用位近似 LRU。}
\exam{常考给页面引用串和页框数计算缺页次数，FIFO 可能有 Belady 异常。}


\section{页面分配与置换策略}
\concept{驻留集与页面分配}
\begin{points}
  \item 驻留集指某进程当前驻留内存的页面集合。
  \item 固定分配：进程运行期间分得的页框数固定。
  \item 可变分配：根据缺页率、工作集等动态调整页框数。
\end{points}
\plain{把这个知识点放回本章主线理解：它回答的是 OS 为了管理资源、隐藏硬件、保证并发正确性而采用的某个对象、规则或步骤。}
\exam{常考选择/填空/简答，让你判断概念归属、比较相近概念，或按“定义-作用-机制-易错点”展开。}


\concept{局部置换与全局置换}
\begin{points}
  \item 局部置换：缺页时只从本进程驻留集中选择页面淘汰。
  \item 全局置换：可从全系统所有可换页面中选择淘汰，进程页框数会动态变化。
  \item 固定分配只能配合局部置换；可变分配可配合局部置换或全局置换。
\end{points}
\plain{内存满了要踢页；OPT 看未来，FIFO 看入队时间，LRU 看最近使用，CLOCK 用引用位近似 LRU。}
\exam{常考给页面引用串和页框数计算缺页次数，FIFO 可能有 Belady 异常。}


\concept{不能组合的策略}
\begin{points}
  \item 固定分配 + 全局置换通常不能组合。
  \item 原因：固定分配要求每个进程页框数不变，而全局置换可能把其他进程页框拿走或让本进程获得额外页框，改变驻留集大小。
\end{points}
\plain{把这个知识点放回本章主线理解：它回答的是 OS 为了管理资源、隐藏硬件、保证并发正确性而采用的某个对象、规则或步骤。}
\exam{常考选择/填空/简答，让你判断概念归属、比较相近概念，或按“定义-作用-机制-易错点”展开。}


\section{工作集、抖动与缺页率}
\concept{工作集模型}
\begin{points}
  \item 工作集是在某段时间窗口内进程实际访问过的页面集合。
  \item 工作集反映进程当前局部性需求，系统应尽量给进程分配足够页框容纳其工作集。
  \item 若分配页框少于工作集需要，缺页率会显著上升。
\end{points}
\plain{给进程的页框太少，它会不停缺页，把 CPU 时间浪费在换页上，这就是抖动。}
\exam{常考抖动原因和解决：降低多道程序度、增加页框、按工作集分配。}


\concept{抖动 thrashing}
\begin{points}
  \item 抖动指系统花费大量时间进行页面换入换出，真正执行用户程序的时间很少。
  \item 现象：CPU 利用率低，磁盘交换区利用率很高，缺页频繁。
  \item 原因：多道程序度过高或每个进程分得页框不足，工作集无法驻留内存。
  \item 解决：降低多道程序度、增加页框、采用工作集策略、改善置换算法。
\end{points}
\plain{给进程的页框太少，它会不停缺页，把 CPU 时间浪费在换页上，这就是抖动。}
\exam{常考抖动原因和解决：降低多道程序度、增加页框、按工作集分配。}


\concept{影响缺页率的因素}
\begin{points}
  \item 页面大小：过小会增加页表和 I/O 次数，过大会增加内部碎片和无用调入。
  \item 分配页框数：一般页框越多缺页率越低，但 FIFO 可能有 Belady 异常。
  \item 页面置换算法：OPT 最优，LRU 通常较好，FIFO 简单但可能异常。
  \item 程序局部性：局部性越好，缺页率越低。
  \item 多道程序度：过高会导致每个进程页框不足，引发抖动。
\end{points}
\plain{页不在内存时才调入；这会产生缺页中断，由 OS 找页、分配页框、必要时置换。}
\exam{常考基本分页与请求分页区别、缺页中断流程、页表项存在位/修改位/引用位。}


